% Preâmbulo — Livro-Texto Dataprev 2026 (Perfil 3 — Desenvolvimento de Software)
% Mesma identidade visual da apostila/ (apostila/preambulo.tex), adaptado para
% material de ensino completo: troca caixas de estratégia de prova (jacaiu,
% peso) por uma caixa de exemplo resolvido.
\usepackage[a4paper,margin=2.2cm,bottom=2.6cm,headheight=14pt]{geometry}
\usepackage[utf8]{inputenc}
\usepackage[T1]{fontenc}
\usepackage[brazilian]{babel}
\usepackage{lmodern}
\usepackage{microtype}
\usepackage{xcolor}
\usepackage{graphicx}
\usepackage{longtable}
\usepackage{array}
\usepackage{colortbl}
\usepackage{booktabs}
\usepackage{multirow}
\usepackage{multicol}
\usepackage{enumitem}
\usepackage{titlesec}
\usepackage{fancyhdr}
\usepackage{needspace}
\usepackage[most]{tcolorbox}
\usepackage{amsmath,amssymb}
\usepackage{tikz}
\usetikzlibrary{shapes,arrows.meta,positioning}
\usepackage[hidelinks,bookmarksnumbered=true,pdfusetitle]{hyperref}
\usepackage{bookmark}

% ---------- Cores (mesma paleta da apostila, + roxo para exemplo) ----------
\definecolor{dpazul}{HTML}{1B3A6B}
\definecolor{dpazulclaro}{HTML}{EAF0FB}
\definecolor{dpverde}{HTML}{1E6B4F}
\definecolor{dpverdeclaro}{HTML}{E7F5EE}
\definecolor{dpvermelho}{HTML}{9B2226}
\definecolor{dpvermelhoclaro}{HTML}{FBE9E9}
\definecolor{dplaranja}{HTML}{B5651D}
\definecolor{dplaranjaclaro}{HTML}{FDF0E3}
\definecolor{dproxo}{HTML}{5B3A8E}
\definecolor{dproxoclaro}{HTML}{EFE9F7}
\definecolor{dpcinza}{HTML}{555555}

% ---------- Cabeçalhos ----------
\pagestyle{fancy}
\fancyhf{}
\renewcommand{\headrulewidth}{0.4pt}
\renewcommand{\chaptermark}[1]{\markboth{\textsc{Capítulo \thechapter}}{}}
\fancyhead[L]{\small\textit{\leftmark}}
\fancyhead[R]{\small Dataprev 2026 --- Teoria completa}
\fancyfoot[C]{\thepage}

\fancypagestyle{plain}{%
  \fancyhf{}
  \fancyfoot[C]{\thepage}
  \renewcommand{\headrulewidth}{0pt}
}

% ---------- Títulos ----------
\titleformat{\chapter}[display]
  {\normalfont\huge\bfseries\color{dpazul}}
  {\filright\Large\textsc{Capítulo \thechapter}}
  {8pt}{\Huge}
\titlespacing*{\chapter}{0pt}{0pt}{28pt}

\titleformat{\section}
  {\normalfont\Large\bfseries\color{dpazul}}{\thesection}{8pt}{}
\titleformat{\subsection}
  {\normalfont\large\bfseries\color{dpazul!80!black}}{\thesubsection}{6pt}{}
\titleformat{\subsubsection}
  {\normalfont\normalsize\bfseries\color{dpcinza}}{\thesubsubsection}{6pt}{}

% ---------- Caixas de destaque ----------
% Ficha do edital (início de capítulo) — igual à apostila, para referência cruzada
\newtcolorbox{edital}{
  enhanced, breakable, colback=dpazulclaro, colframe=dpazul,
  boxrule=0.8pt, arc=2pt, left=8pt, right=8pt, top=6pt, bottom=6pt,
  fonttitle=\bfseries, coltitle=white, colbacktitle=dpazul,
  title={Ficha do edital},
  before skip=10pt, after skip=12pt
}
% Conceito-chave — a caixa principal deste livro, usada extensivamente
\newtcolorbox{conceito}[1][]{
  enhanced, breakable, colback=white, colframe=dpazul!70!black,
  boxrule=0.6pt, arc=2pt, left=8pt, right=8pt, top=5pt, bottom=5pt,
  fonttitle=\bfseries\color{white}, coltitle=white,
  colbacktitle=dpazul, title={#1},
  before skip=8pt, after skip=10pt
}
% Cuidado / confusão comum (mesma caixa "pegadinha" da apostila, retitulada:
% aqui é erro conceitual comum, não tática de banca)
\newtcolorbox{cuidado}[1][Cuidado --- confusão comum]{
  enhanced, breakable, colback=dpvermelhoclaro, colframe=dpvermelho,
  boxrule=0.7pt, arc=2pt, left=8pt, right=8pt, top=6pt, bottom=6pt,
  fonttitle=\bfseries, coltitle=white, colbacktitle=dpvermelho,
  title={#1},
  before skip=10pt, after skip=10pt
}
% Regra prática / como fixar
\newtcolorbox{regrapratica}[1][Regra prática]{
  enhanced, breakable, colback=dpverdeclaro, colframe=dpverde,
  boxrule=0.7pt, arc=2pt, left=8pt, right=8pt, top=6pt, bottom=6pt,
  fonttitle=\bfseries, coltitle=white, colbacktitle=dpverde,
  title={#1},
  before skip=10pt, after skip=10pt
}
% Exemplo resolvido
\newtcolorbox{exemplo}[1][Exemplo]{
  enhanced, breakable, colback=dproxoclaro, colframe=dproxo,
  boxrule=0.7pt, arc=2pt, left=8pt, right=8pt, top=6pt, bottom=6pt,
  fonttitle=\bfseries, coltitle=white, colbacktitle=dproxo,
  title={#1},
  before skip=10pt, after skip=10pt
}

% ---------- Tabelas: cabeçalho colorido ----------
\arrayrulecolor{dpazul!40}
\renewcommand{\arraystretch}{1.25}

% ---------- Listas ----------
\setlist{topsep=4pt,itemsep=2pt,parsep=0pt}

% ---------- Parágrafos ----------
\setlength{\parindent}{0pt}
\setlength{\parskip}{4pt}

% ---------- Sumário: profundidade ----------
\setcounter{tocdepth}{2}
\setcounter{secnumdepth}{3}
